\chapter{Identity Axis, Normalized Cycles, and the Typed Relation Solver}
\label{ch:identity-solver}

\section{Purpose}
The preceding chapters establish several mathematically distinct routes that meet at the same balanced coordinate.  This chapter does not add a new proof of the Riemann hypothesis.  It does two narrower jobs.  First, it separates the \emph{identity coordinate} from phase evolution by introducing a normalized recurrence coordinate before radians.  Second, it formalizes the implication structure as a typed solver so that exact theorems, standard structures, model assignments, and open bridges cannot be silently composed as if they had the same epistemic status.

The corresponding executable is
\begin{center}
\path{TIR/zeta_information_axis/src/critical_axis/solver.py}.
\end{center}
The runtime counterpart is the CIEL/PhaseNav geometry-dependency layer in the companion \texttt{noema-phasenav-core} repository.  The two layers have deliberately different responsibilities: TIR records what follows from what; PhaseNav supplies geometric states, tangent vectors, normalized cycles, relation paths, and holonomy candidates.

\section{The identity axis is not a frozen state}
Let the binary identity coordinate be $\Sigma$.  The distinguished balance axis is
\begin{equation}
\Sigma_\star=\frac12.
\label{eq:sigma-axis}
\end{equation}
The meaning of ``axis'' is crucial.  A system may move along an orbit while preserving the level-set condition that identifies the relevant sector.  In the qubit reduction, fixing $\sigma=1/2$ fixes
\begin{equation}
n_z=2\sigma-1=0,
\end{equation}
while the azimuthal phase remains free.  Thus identity need not mean a static state.  It means that phase evolution occurs on an invariant leaf selected by the identity coordinate.

Introduce the centered coordinate
\begin{equation}
x=\sigma-\frac12.
\label{eq:centered-sigma}
\end{equation}
Then binary complement becomes reflection,
\begin{equation}
\sigma\mapsto1-\sigma
\qquad\Longleftrightarrow\qquad
x\mapsto-x.
\label{eq:sigma-reflection}
\end{equation}
This gives a primitive geometric interpretation of sign: $x<0$ and $x>0$ are opposite orientations from the identity axis, while $x=0$ is the fixed axis itself.  In the PhaseNav dependency package this principle is implemented by \texttt{IdentityAxis} and \texttt{OrientedScalar} rather than by assigning semantic meaning to a sign after the fact.

\section{Frequency before angle}
Let
\begin{equation}
q\in\mathbb R/\mathbb Z
\label{eq:q-cycle}
\end{equation}
be a normalized cycle coordinate.  One recurrence is $q\mapsto q+1$.  The intrinsic frequency is therefore the winding rate
\begin{equation}
f_q=\frac{\Delta N}{\Delta t},
\label{eq:winding-frequency}
\end{equation}
which requires neither radians nor a prior numerical value of $\pi$.

Only after the recurrent topology has been defined do we choose an angular representation
\begin{equation}
\phi=Cq,
\label{eq:closure-C}
\end{equation}
where $C$ is the closure period in the chosen angular unit.  Standard radians correspond to $C=2\pi$.  This ordering is conceptually useful because it distinguishes three layers:
\begin{equation}
\boxed{
\text{recurrence}\;\longrightarrow\;\text{winding}\;\longrightarrow\;\text{angular representation}
}.
\end{equation}
The topology of a cycle is therefore prior to its radian parametrization.

\section{Spinorial double cover}
For a spin-$1/2$ lift, one projective cycle changes the spinor sheet while two projective cycles restore it:
\begin{equation}
+1\longrightarrow-1\longrightarrow+1.
\label{eq:spinor-sheet-cycle}
\end{equation}
Equivalently,
\begin{equation}
T_{\rm spin}=2T_{\rm proj}.
\label{eq:double-period}
\end{equation}
This statement can be made before introducing $2\pi$ and $4\pi$.  The latter are the radian representatives obtained only after applying \eqref{eq:closure-C} with $C=2\pi$.

\section{Independent routes to the half axis}
Within their stated domains, four routes select the same value:
\begin{enumerate}
\item \textbf{Complement fixedness:}
\begin{equation}
\sigma=1-\sigma\quad\Longrightarrow\quad\sigma=\frac12.
\end{equation}
\item \textbf{Shannon stationarity and concavity:}
\begin{equation}
H'_{\rm bin}(\sigma)=0,
\qquad H''_{\rm bin}(\sigma)<0
\quad\Longrightarrow\quad
\sigma=\frac12.
\end{equation}
\item \textbf{Equatorial Berry holonomy:} for the state family of Chapter~\ref{ch:bloch},
\begin{equation}
\mathcal U_B(\sigma)=-1
\quad\Longrightarrow\quad
\sigma=\frac12
\end{equation}
for the nontrivial interior branch.
\item \textbf{Symmetric complementary cancellation:}
\begin{equation}
\sqrt\sigma-e^{0}\sqrt{1-\sigma}=0
\quad\Longrightarrow\quad
\sigma=\frac12,
\end{equation}
where the minus sign is the relative half-turn phase of Theorem~\ref{thm:cancellation}.
\end{enumerate}
The executable solver evaluates these routes independently and verifies their common solution.  Their agreement is a consistency result inside the declared representations; it is not permission to bypass the open zeta zero-state bridge.

\section{Arithmetic--geometry factorization of the information normalization}
Write the information assigned to one normalized projective turn as
\begin{equation}
\frac{\dd I}{\dd q}=\frac{\ln2}{12}.
\label{eq:info-per-turn}
\end{equation}
Equation~\eqref{eq:info-per-turn} is a TIR/Metatime model assignment, not a theorem of Shannon theory.  Given that assignment and the angular closure \eqref{eq:closure-C}, the density with respect to angular phase is exactly
\begin{equation}
\frac{\dd I}{\dd\phi}
=
\frac{\ln2}{12C}.
\label{eq:kappa-C}
\end{equation}
For the radian representation $C=2\pi$ this becomes
\begin{equation}
\boxed{
\kappa=\frac{\ln2}{24\pi}
}.
\label{eq:kappa-cross-derived}
\end{equation}
Thus the arithmetic and geometric factors are explicitly separated:
\begin{equation}
\boxed{
24\pi
=
\underbrace{12}_{\text{declared information-cycle count}}
\underbrace{2\pi}_{\text{geometric closure in radians}}.
}
\end{equation}

The integer factor has several exact factorizations,
\begin{equation}
24=8\cdot3=12\cdot2=6\cdot4.
\label{eq:24-cross-factorization}
\end{equation}
The equality of the integers is exact arithmetic.  Assigning the factors respectively to eight mixing sectors and three flavours, twelve projective cycles and two half-turns, or six spinor cycles and four half-turns is model semantics.  The solver therefore records those assignments as \status{MODEL}, not as independent theorem-level derivations.

\section{Paired implication graph}
The relevant $1/2$ slice is better represented as a graph than as a list of isolated coincidences.  The principal exact or standard edges are
\begin{align}
\sigma=\frac12
&\Longrightarrow C(\sigma)=\sigma,\\
\sigma=\frac12
&\Longrightarrow H_{\rm bin}(\sigma)=\ln2,\\
\sigma=\frac12+\text{qubit representation}
&\Longrightarrow n_z=0,\\
n_z=0+\text{one equatorial loop}
&\Longrightarrow \mathcal U_B=-1,\\
\sigma=\frac12+\text{symmetric readout}+\text{half-turn}
&\Longrightarrow \mathcal A=0,\\
\text{centred zeta chart}+\mathcal I(u)=1/u
&\Longrightarrow \text{critical-axis set invariance},\\
\text{projective cycle}+\text{spin }\frac12
&\Longrightarrow \text{double cover}.
\end{align}
The graph may contain several edges incident on the same node.  ``Everything has a pair'' therefore means that every declared concept participates in at least one typed relational edge; it does \emph{not} mean that every concept has exactly one exclusive partner.

\section{Holonomy is an edge property, not a universal decoration}
A semantic relation is not automatically holonomic.  A holonomy is attached only when a connection and a closed or composable path are defined.  This prevents category errors such as assigning an arbitrary $\pi$ phase to an energy--matter pair merely because both are part of the same ontology.

For a path $P=(v_0,\ldots,v_n)$ whose edges carry normalized phase increments $h_k\in\mathbb R/\mathbb Z$, the composable holonomy coordinate is
\begin{equation}
h(P)=\sum_{k=0}^{n-1}h_k\pmod1.
\label{eq:path-holonomy-turns}
\end{equation}
A complex $U(1)$ representative may be chosen afterwards.  The PhaseNav dependency layer stores the normalized turn first; this keeps the primitive relation independent of radian convention.

\section{Typed closure and the RH firewall}
Each solver rule has a status in
\begin{equation}
\{\status{EXACT},\status{STANDARD},\status{MODEL},\status{OPEN}\}.
\end{equation}
The default closure engine permits only \status{EXACT} and \status{STANDARD} rules.  Model closure additionally permits \status{MODEL}.  \status{OPEN} rules are never automatically promoted.

The critical firewall is therefore explicit.  The chain
\begin{equation}
\text{zero-state representation}
+\text{exact cancellation}
+\sigma=\Ree s
\Longrightarrow
\text{all nontrivial zeros on }\Ree s=\frac12
\end{equation}
is stored as \status{OPEN}.  Consequently neither exact nor model closure can derive the Riemann hypothesis.  Instead the solver reports \texttt{zero\_state\_representation} as the missing premise.

\section{TIR crosswalk}
The parent Metatime/TIR monograph supplies the broader structures used here:
\begin{itemize}
\item the information normalization $\kappa$ and its publication claim boundary;
\item the Poincar\'e-disk geometry and phase-holonomy language;
\item three-flavour/neutrino structures;
\item sector-holonomy operators and their retrospective/prospective status discipline.
\end{itemize}
The parent monograph must therefore be read together with the dedicated critical-axis monograph.  The direction of implication is asymmetric: the parent TIR supplies model-level structural assignments, while the present monograph supplies exact half-axis theorems and returns only those consequences that satisfy the solver's status gate.

\section{Implementation contract}
The executable contract is:
\begin{enumerate}
\item geometry creates centered and tangent coordinates;
\item recurrence creates $q\in\mathbb R/\mathbb Z$ and winding frequency;
\item a representation map may convert turns to radians;
\item vectorization consumes geometric tangent coordinates rather than opaque semantic hashes;
\item relation edges carry explicit status and optional holonomy;
\item the solver composes only status-authorized hyperedges;
\item generated PhaseNav paths remain candidates until a TIR rule supplies independent semantic provenance.
\end{enumerate}
This gives a concrete separation between geometry, semantics, and proof status while preserving the exploratory role of the broader TIR programme.
